\documentclass[11pt]{article}

\usepackage{amsmath,amssymb,amsthm,mathtools}
\usepackage[margin=1in]{geometry}
\usepackage{microtype}
\usepackage[hidelinks]{hyperref}

\newtheorem{theorem}{Theorem}[section]
\newtheorem{lemma}[theorem]{Lemma}
\newtheorem{corollary}[theorem]{Corollary}

\newcommand{\dd}{\,\mathrm{d}}
\newcommand{\one}{\mathbf{1}}
\newcommand{\ip}[2]{\left\langle #1,#2\right\rangle}
\newcommand{\join}{\mathbin{\vee}}

\title{Six Explicit Verified-Base Cone Cases}
\author{}
\date{}

\begin{document}
\maketitle

\begin{abstract}
We prove a conditional cone lemma from first principles and apply it to
six explicitly verified five-vertex bases.  The bases are Atlas 33, 34,
36, 40, 45, and 49.  Each already has an accepted graphon lower bound,
and its exact chromatic target is explicitly factorized and checked to
be nonnegative, increasing, and convex from its chromatic threshold.
The result proves the previously open cone graphs Atlas 177, 179, 183,
192, 200, and 205.  No general universal-vertex closure is asserted.
\end{abstract}


%%%%%%%%%%%%%%%%%%%%%%%%%%%%%%%%%%%%%%%%%%%%%%%%%%%%%%%%%%%%%%%%%%%%%%%%%%%%%%%
\section{A conditional cone lemma}
%%%%%%%%%%%%%%%%%%%%%%%%%%%%%%%%%%%%%%%%%%%%%%%%%%%%%%%%%%%%%%%%%%%%%%%%%%%%%%%

Let $W\colon\Omega^2\to[0,1]$ be a symmetric measurable graphon on a
probability space $(\Omega,\mu)$.  For a finite graph $G$, write
\[
t(G,W)=
\int_{\Omega^{V(G)}}
\prod_{uv\in E(G)}W(x_u,x_v)
\prod_{u\in V(G)}\dd\mu(x_u).
\]
Put
\[
p=t(K_2,W),
\qquad
d(x)=\int_\Omega W(x,y)\dd\mu(y).
\]
Define
\[
\tau(x)=
\int_{\Omega^2}W(x,y)W(x,z)W(y,z)
\dd\mu(y)\dd\mu(z).
\]

\begin{lemma}[Weighted rooted-triangle inequality]
\label{lem:weighted}
If $p>0$, then for every integer $h\geq0$,
\[
\int_\Omega d(x)^h\tau(x)\dd\mu(x)
\geq
\frac{2p-1}{p}
\int_\Omega d(x)^{h+2}\dd\mu(x).
\]
\end{lemma}

\begin{proof}
Put $M_j=\int d^j$, let $T_W$ be the graphon operator, and write
\[
I_h=\int d^h\tau,
\qquad
J_h=\ip{d^h}{T_Wd}.
\]
Applying $ab\geq a+b-1$ to $W(x,y),W(x,z)$, multiplying by
$d(x)^hW(y,z)$, and integrating gives
\begin{equation}
I_h\geq2J_h-pM_h.
\label{eq:I-J}
\end{equation}
Set $U=1-W$, $u=T_U\one=1-d$, and $q=1-p$.  Then
\[
T_Wd=d-q\one+T_Uu,
\]
so
\begin{equation}
J_h=M_{h+1}-qM_h+\ip{d^h}{T_Uu}.
\label{eq:J}
\end{equation}
By symmetry of $U$, twice
\[
\ip{d^h}{T_Uu}-\int d^hu^2
\]
equals
\[
\int_{\Omega^2}U(x,y)
\bigl(u(y)-u(x)\bigr)
\bigl(d(x)^h-d(y)^h\bigr)
\dd\mu(x)\dd\mu(y)\geq0,
\]
because $d=1-u$.  Therefore \eqref{eq:J} gives
\[
J_h\geq M_{h+2}-M_{h+1}+pM_h.
\]
Using this in \eqref{eq:I-J}, and then subtracting
$(2p-1)M_{h+2}/p$, leaves exactly
\[
\frac1p\int_\Omega d(x)^h(d(x)-p)^2\dd\mu(x)\geq0.
\]
\end{proof}

For $d(x)>0$, define
\[
\dd\mu_x(y)=\frac{W(x,y)}{d(x)}\dd\mu(y),
\]
and let $W_x$ denote $W$ on $(\Omega,\mu_x)$.  Its edge density is
$z_x=\tau(x)/d(x)^2$.  Put $z_x=0$ when $d(x)=0$.

\begin{corollary}[Weighted link average]
\label{cor:average}
For every integer $N\geq2$,
\[
\int_\Omega d(x)^Nz_x\dd\mu(x)
\geq
\left(2-\frac1p\right)
\int_\Omega d(x)^N\dd\mu(x).
\]
\end{corollary}

\begin{proof}
Apply Lemma~\ref{lem:weighted} with $h=N-2$.
\end{proof}

\begin{lemma}[Verified-base cone lemma]
\label{lem:cone}
Let $F$ be a graph on $N\geq2$ vertices and let
\[
\phi(z)=\Phi_F(z).
\]
Assume that for some $a\in[0,1)$:
\begin{enumerate}
\item every graphon $V$ of edge density $z\geq a$ satisfies
$t(F,V)\geq\phi(z)$;
\item on $[a,1]$, one has $\phi\geq0$, $\phi'\geq0$, and
$\phi''\geq0$;
\item $\phi(a)=0$.
\end{enumerate}
Then every graphon $W$ of edge density
\[
p\geq b:=\frac1{2-a}
\]
satisfies
\[
t(K_1\join F,W)\geq\Phi_{K_1\join F}(p).
\]
\end{lemma}

\begin{proof}
Put $c=2-1/p$.  The assumption $p\geq1/(2-a)$ is exactly $c\geq a$.
Convexity shows that every graphon $V$ of arbitrary edge density $z$
satisfies the tangent bound
\[
t(F,V)\geq\phi(c)+\phi'(c)(z-c).
\]
Indeed, for $z\geq a$ this follows from the assumed graphon bound and
convexity.  For $z<a$, the tangent is nonpositive because it is at most
$\phi(a)=0$ at $a$ and has nonnegative slope, whereas $t(F,V)\geq0$.

Conditioning on the image $x$ of the new cone vertex gives
\[
t(K_1\join F,W)=\int_\Omega d(x)^N t(F,W_x)\dd\mu(x).
\]
Let $M_N=\int d^N$.  Applying the tangent bound in every link,
\begin{align*}
t(K_1\join F,W)
&\geq M_N\phi(c)
\\
&\quad+
\phi'(c)\left(\int d(x)^Nz_x\dd\mu(x)-cM_N\right).
\end{align*}
The second line is nonnegative by Corollary~\ref{cor:average}.  Jensen
then gives
\[
t(K_1\join F,W)\geq p^N\phi(c).
\]

If $q=1-p$, then $1-c=q/p$.  Since
$\chi_{K_1\join F}(x)=x\chi_F(x-1)$,
\begin{align*}
\Phi_{K_1\join F}(p)
&=q^N\chi_F\!\left(\frac pq\right)
\\
&=p^N\left(\frac qp\right)^N
\chi_F\!\left(\frac p q\right)
=p^N\phi(c).
\end{align*}
\end{proof}


%%%%%%%%%%%%%%%%%%%%%%%%%%%%%%%%%%%%%%%%%%%%%%%%%%%%%%%%%%%%%%%%%%%%%%%%%%%%%%%
\section{The six verified bases}
%%%%%%%%%%%%%%%%%%%%%%%%%%%%%%%%%%%%%%%%%%%%%%%%%%%%%%%%%%%%%%%%%%%%%%%%%%%%%%%

Let $F_j$ denote NetworkX Atlas graph $j$.  The following table gives
the complete input verification.  The graphon bound column records the
accepted reason already present in the catalogue before this theorem.

\begin{center}
\begin{tabular}{cclll}
\hline
$j$ & graph6 & accepted bound & $a$ & $\phi_j(z)=\Phi_{F_j}(z)$ \\
\hline
33 & \texttt{Dn?} & diamond plus isolate product & $1/2$ & $z(2z-1)^2$ \\
34 & \texttt{D@\char123} & rooted common leaves & $1/2$ & $z^3(2z-1)$ \\
36 & \texttt{DJc} & rooted two-edge tail & $1/2$ & $z^3(2z-1)$ \\
40 & \texttt{DjW} & forest-cone theorem & $1/2$ & $z^2(2z-1)^2$ \\
45 & \texttt{DJ\char123} & clique common leaf & $2/3$ & $z^2(2z-1)(3z-2)$ \\
49 & \texttt{Df\char123} & special paw cone & $2/3$ & $z(2z-1)^2(3z-2)$ \\
\hline
\end{tabular}
\end{center}

The chromatic polynomials, obtained either from the displayed structures
or by deletion--contraction, are
\begin{align*}
\chi_{F_{33}}(x)&=x^2(x-1)(x-2)^2,\\
\chi_{F_{34}}(x)=\chi_{F_{36}}(x)&=x(x-1)^3(x-2),\\
\chi_{F_{40}}(x)&=x(x-1)^2(x-2)^2,\\
\chi_{F_{45}}(x)&=x(x-1)^2(x-2)(x-3),\\
\chi_{F_{49}}(x)&=x(x-1)(x-2)^2(x-3).
\end{align*}
Substitution in the definition of $\Phi$ gives exactly the table.

For completeness, the required shape conditions are also checked
explicitly.  For the first four rows, every displayed factor is
nonnegative and nondecreasing on $[1/2,1]$; for the last two, on
$[2/3,1]$.  A product of nonnegative nondecreasing affine factors is
nonnegative, nondecreasing, and convex: its first and second derivatives
are sums of products of nonnegative factors and positive constant
slopes.  Repeated factors cause no problem.  Each polynomial vanishes
at its displayed $a$.  Thus every hypothesis of
Lemma~\ref{lem:cone} is verified, not inferred from tests.


%%%%%%%%%%%%%%%%%%%%%%%%%%%%%%%%%%%%%%%%%%%%%%%%%%%%%%%%%%%%%%%%%%%%%%%%%%%%%%%
\section{New positive cases}
%%%%%%%%%%%%%%%%%%%%%%%%%%%%%%%%%%%%%%%%%%%%%%%%%%%%%%%%%%%%%%%%%%%%%%%%%%%%%%%

\begin{theorem}[Six explicit cones]
\label{thm:six}
For
\[
j\in\{33,34,36,40,45,49\},
\]
the cone $K_1\join F_j$ satisfies
\[
t(K_1\join F_j,W)\geq\Phi_{K_1\join F_j}(p)
\]
for every graphon throughout its full required density interval.
\end{theorem}

\begin{proof}
For $j=33,34,36,40$, the verified-base threshold is $a=1/2$, so
Lemma~\ref{lem:cone} gives the cone threshold
\[
b=\frac1{2-1/2}=\frac23.
\]
Each base is $3$-chromatic, so each cone is $4$-chromatic and this is
the exact required threshold.

For $j=45,49$, one has $a=2/3$, and the cone threshold is
\[
b=\frac1{2-2/3}=\frac34.
\]
Each base is $4$-chromatic, so each cone is $5$-chromatic and this is
again the exact required threshold.
\end{proof}

\begin{corollary}[Catalogue IDs]
The newly positive Atlas graphs are
\[
\begin{array}{c|c|c|c}
\text{Atlas ID}&\text{graph6}&\text{base Atlas}&\text{threshold}\\
\hline
177&\texttt{En\char123 G}&33&2/3\\
179&\texttt{E\_\char126 w}&34&2/3\\
183&\texttt{EjvG}&36&2/3\\
192&\texttt{E\char126 Xo}&40&2/3\\
200&\texttt{E\char126\char123 W}&45&3/4\\
205&\texttt{E\char126\char126 G}&49&3/4
\end{array}
\]
\end{corollary}

At the threshold the target vanishes, and the balanced complete
$(\chi(H)-1)$-partite graphon attains equality.  At $p=1$, both sides
equal $1$.  At every balanced multipartite grid point in between,
equality follows from the chromatic-polynomial identity.


%%%%%%%%%%%%%%%%%%%%%%%%%%%%%%%%%%%%%%%%%%%%%%%%%%%%%%%%%%%%%%%%%%%%%%%%%%%%%%%
\section{Scope and hidden assumptions}
%%%%%%%%%%%%%%%%%%%%%%%%%%%%%%%%%%%%%%%%%%%%%%%%%%%%%%%%%%%%%%%%%%%%%%%%%%%%%%%

The proof applies Lemma~\ref{lem:cone} only to the six bases whose exact
graphon inequalities and analytic target properties were verified in
the table.  It does not treat arbitrary positive bases and does not
assert a general universal-vertex or join closure.

Individual graphon links may have density below the base threshold; the
tangent proof handles those links by a nonpositive lower bound.  The
argument uses ordinary homomorphism densities directly and assumes no
regularity, step structure, injectivity, minimum degree, or pointwise
degree lower bound.

%%%%%%%%%%%%%%%%%%%%%%%%%%%%%%%%%%%%%%%%%%%%%%%%%%%%%%%%%%%%%%%%%%%%%%%%%%%%%%%
\appendix
\section{What the Lean formalization actually proves}
%%%%%%%%%%%%%%%%%%%%%%%%%%%%%%%%%%%%%%%%%%%%%%%%%%%%%%%%%%%%%%%%%%%%%%%%%%%%%%%

Theorem~\ref{thm:six} is formalized for all six cones; Atlas $177$, $179$,
$183$, $192$, $200$ and $205$ are \texttt{verified}.  The development is
\texttt{lean/Taeyoung/Methods/BaseCone/}, over the abstract
\texttt{Methods/ConeBound.lean}; Atlas $192$ is in
\texttt{Methods/ForestCone/Rows.lean} instead, because its base bound is the
forest-cone theorem.  One difference of substance.

\subsection*{Hypothesis (2) of Lemma~\ref{lem:cone} is replaced, not verified}

The convexity hypothesis $\phi\geq0$, $\phi'\geq0$, $\phi''\geq0$ on $[a,1]$ is
not checked for the six targets.  It is replaced by an algebraic one: each
$\Phi_{F_j}$ is a product of affine factors
\[
 \phi(z)=\prod_i\bigl(1-k_i(1-z)\bigr),\qquad k_i\geq0,
\]
and the cone lemma is stated in Lean for that class
(\texttt{satisfiesLowerBound\_of\_baseCone}).  The six applications then reduce
to naming a list of slopes $k_i$, with no second derivative and no convexity
theory anywhere in the development.

The affine form also improves the last step.  Lemma~\ref{lem:cone} bounds
$p^{N}\phi(2-1/p)$ below by $\Phi_{K_1\join F}(p)$; in the affine form the two
are \emph{equal}, which is \texttt{pow\_mul\_affineProd\_shift}, so the cone's
own target is reached exactly rather than through an estimate.

One base needs no separate input: Atlas $33$ is the diamond together with an
isolated vertex, which is itself the cone over $K_{1,2}$ with target $z^2$, so
the same theorem supplies its bound.

\end{document}
